\section{Proposed Solution}

\subsection{Design Overview}
Fig.~\ref{fig:datapath} shows the processor. It is a Harvard machine: instruction memory
and data memory are separate 1024-word arrays, indexed by \texttt{addr[11:2]}, with
combinational reads. All registers update on the rising clock edge and use a
synchronous, active-high reset. The supported subset is listed in Table~\ref{tab:isa}.

\begin{table}[!t]
\centering
\caption{Supported RV32I Instruction Subset (29 Instructions)}
\label{tab:isa}
\renewcommand{\arraystretch}{1.1}
\begin{tabular}{@{}p{1.5cm}p{6.2cm}@{}}
\toprule
\textbf{Type} & \textbf{Instructions} \\ \midrule
R-type   & ADD, SUB, AND, OR, XOR, SLL, SRL, SRA, SLT, SLTU \\
I-type   & ADDI, ANDI, ORI, XORI, SLTI, SLTIU, SLLI, SRLI, SRAI \\
Load/Store & LW, SW (word-aligned) \\
Branch   & BEQ, BNE, BLT, BGE \\
Jump     & JAL, JALR \\
U-type   & LUI \\ \bottomrule
\end{tabular}
\end{table}

% ---------------- Datapath figure ----------------
\begin{figure*}[!t]
\centering
\begin{tikzpicture}[
  font=\scriptsize, >={Stealth[length=1.5mm]},
  blk/.style={draw, rounded corners=2pt, minimum width=1.45cm, minimum height=0.7cm, align=center, fill=white},
  preg/.style={draw, fill=gray!30, minimum width=0.36cm, minimum height=4.6cm},
  mux/.style={draw, fill=white, trapezium, trapezium angle=72, shape border rotate=270, minimum height=0.95cm, inner sep=1pt},
  w/.style={->, thick},
  ctl/.style={->, dashed, red!70!black},
  fA/.style={->, thick, blue!75!black},
  fB/.style={->, thick, teal!70!black},
  lbl/.style={font=\tiny, inner sep=1pt}
]
% stage backgrounds
\fill[stIF]  (-0.65,-3.85) rectangle (2.75,2.65);
\fill[stID]  (3.25,-3.85) rectangle (5.95,2.65);
\fill[stEX]  (6.45,-3.85) rectangle (9.55,2.65);
\fill[stMEM] (10.05,-3.85) rectangle (12.15,2.65);
\fill[stWB]  (12.65,-3.85) rectangle (14.35,2.65);
\foreach \x/\t in {1.05/IF, 4.6/ID, 8.0/EX, 11.1/MEM, 13.5/WB}
  \node[font=\footnotesize\bfseries] at (\x,2.9) {\t};
% pipeline registers
\node[preg] (r1) at (3.0,-0.15)  {}; \node[rotate=90] at (3.0,-0.15) {IF/ID};
\node[preg] (r2) at (6.2,-0.15)  {}; \node[rotate=90] at (6.2,-0.15) {ID/EX};
\node[preg] (r3) at (9.8,-0.15)  {}; \node[rotate=90] at (9.8,-0.15) {EX/MEM};
\node[preg] (r4) at (12.4,-0.15) {}; \node[rotate=90] at (12.4,-0.15) {MEM/WB};
% IF
\node[mux] (pcm) at (-0.25,0) {};
\node[blk, minimum width=0.7cm] (pc) at (0.6,0) {PC};
\node[blk] (im) at (1.85,0) {Instr.\\Memory};
\node[blk, minimum width=0.75cm, minimum height=0.45cm] (p4) at (1.85,1.3) {+4};
\draw[w] (pcm.east) -- (pc.west);
\draw[w] (pc.east) -- (im.west);
\draw[w] (pc.east) ++(0.12,0) |- (p4.west);
\fill (pc.east) ++(0.12,0) circle (1.2pt);
\draw[w] (p4.east) -- (2.35,1.3) -- (2.35,1.9) -- (-0.45,1.9) -- (-0.45,0.25) -- ($(pcm.west)+(0,0.25)$);
\node[lbl] at (0.6,2.02) {pc+4};
\draw[w] (im.east) -- (r1.west |- im.east);
% ID
\node[blk] (cu) at (4.6,1.3) {Control\\Unit};
\node[blk, minimum height=1.05cm] (rf) at (4.6,0) {Register File\\(write-first\\bypass)};
\node[blk] (ig) at (4.6,-1.35) {Imm.\\Generator};
\draw[w] (r1.east |- cu.west) -- (cu.west);
\draw[w] (r1.east |- 0,0.2) -- (rf.west |- 0,0.2);
\draw[w] (r1.east |- ig.west) -- (ig.west);
\draw[w] (cu.east) -- (r2.west |- cu.east);
\draw[w] ($(rf.east)+(0,0.25)$) -- node[lbl,above]{rd1} (r2.west |- 0,0.25);
\draw[w] ($(rf.east)+(0,-0.25)$) -- node[lbl,below]{rd2} (r2.west |- 0,-0.25);
\draw[w] (ig.east) -- node[lbl,above]{imm} (r2.west |- ig.east);
% control pass-through
\draw[w, gray] (r2.east |- 0,2.05) -- node[lbl,above,text=black]{control} (r3.west |- 0,2.05);
\draw[w, gray] (r3.east |- 0,2.05) -- (r4.west |- 0,2.05);
% EX
\node[mux] (ma) at (7.35,0.8) {};
\node[mux, minimum height=1.25cm] (mb) at (7.35,-1.0) {};
\node[blk, minimum width=0.95cm, minimum height=2.0cm] (alu) at (8.65,-0.1) {ALU};
\node[blk, minimum width=1.25cm, minimum height=0.55cm] (bu) at (8.65,1.35) {Branch Unit};
\draw[w] (r2.east |- 0,1.05) -- (ma.west |- 0,1.05);
\draw[w] (r2.east |- 0,-0.6) -- (mb.west |- 0,-0.6);
\draw[w] (ma.east) -- (alu.west |- ma.east);
\draw[w] (mb.east) -- (alu.west |- mb.east);
\draw[w] (r2.east |- 0,-1.42) -- node[lbl,below,pos=0.3]{imm} (mb.west |- 0,-1.42);
\draw[w] (ma.east) ++(0.15,0) |- (bu.west);
\fill (ma.east) ++(0.15,0) circle (1.2pt);
\draw[w] (alu.east) -- node[lbl,above]{result} (r3.west |- alu.east);
\draw[w] (mb.east) ++(0.15,0) -- ++(0,-0.35) -- (r3.west |- 0,-1.35);
\fill (mb.east) ++(0.15,0) circle (1.2pt);
\node[lbl] at (9.3,-1.22) {store};
% MEM
\node[blk, minimum height=1.2cm] (dm) at (11.1,0) {Data\\Memory};
\draw[w] (r3.east |- 0,0.2) -- node[lbl,above]{addr} (dm.west |- 0,0.2);
\draw[w] (r3.east |- 0,-0.3) -- node[lbl,below]{wdata} (dm.west |- 0,-0.3);
\draw[w] (dm.east) -- node[lbl,above]{rdata} (r4.west |- dm.east);
\draw[w] (r3.east |- 0,-1.2) -- (r4.west |- 0,-1.2);
\node[lbl] at (11.1,-1.05) {ALU result};
% WB
\node[mux] (wm) at (13.3,-0.3) {};
\draw[w] (r4.east |- 0,0) -- (wm.west |- 0,0);
\draw[w] (r4.east |- 0,-1.2) -- ++(0.3,0) |- (wm.west |- 0,-0.6);
% write-back path
\draw[w] (wm.east) -- (13.95,-0.3) -- (13.95,-2.45) -- (3.45,-2.45) -- (3.45,-0.35) -- (rf.west |- 0,-0.35);
\node[lbl, fill=stWB] at (13.55,-2.3) {wb\_data};
% branch / jump redirect
\draw[w, red!70!black] (bu.north) -- (8.65,2.4) -- (-0.6,2.4) -- (-0.6,-0.25) -- ($(pcm.west)+(0,-0.25)$);
\node[lbl, text=red!70!black, fill=stID] at (4.6,2.4) {branch\_target, pc\_src};
% forwarding paths
\fill[blue!75!black] (r3.east |- 0,-1.2) ++(0.2,0) circle (1.3pt);
\draw[fA] (r3.east |- 0,-1.2) ++(0.2,0) -- ++(0,-0.75) -- (6.6,-1.95) -- (6.6,0.75) -- (ma.west |- 0,0.75);
\draw[fA] (6.6,-0.88) -- (mb.west |- 0,-0.88);
\node[lbl, text=blue!75!black, fill=stEX] at (8.3,-2.07) {EX/MEM forward (mem\_fwd\_data)};
\fill[teal!70!black] (6.8,-2.45) circle (1.3pt);
\draw[fB] (6.8,-2.45) -- (6.8,0.55) -- (ma.west |- 0,0.55);
\draw[fB] (6.8,-1.15) -- (mb.west |- 0,-1.15);
\node[lbl, text=teal!70!black, fill=stMEM] at (11.1,-2.6) {MEM/WB forward (wb\_data)};
% hazard and forwarding units
\node[blk, fill=hz, minimum width=1.7cm] (fu) at (8.3,-3.3) {Forwarding Unit};
\node[blk, fill=hz, minimum width=1.9cm] (hu) at (4.6,-3.3) {Hazard Detection Unit};
\draw[ctl] (7.6,-2.95) -- node[lbl, right, pos=0.35]{forward\_a/b} (7.6,-1.6);
\draw[ctl] (hu.west) -| node[lbl, pos=0.75, right]{pc\_write} (pc.south);
\draw[ctl] (hu.west) ++(0,0.15) -| node[lbl, pos=0.85, left]{if\_id\_write, flush} (r1.south);
\draw[ctl] (hu.east) -| node[lbl, pos=0.85, right]{flush} (r2.south);
\end{tikzpicture}
\caption{Block diagram of the five-stage pipelined RV32I processor. Grey bars are pipeline
registers and yellow blocks are the hazard-handling units. Blue and teal lines are the
EX/MEM and MEM/WB forwarding paths, the red line is the branch/jump redirect, and dashed
lines are hazard control signals.}
\label{fig:datapath}
\end{figure*}

\subsection{Design Methodology}
The work followed an \emph{interface-first} flow (Fig.~\ref{fig:flow}). Before any
module was written, the team fixed a shared header (\texttt{defines.vh}) and agreed on
the Verilog ports of every module, the control truth table and the hazard policies.
Each member then built and unit-tested their modules on a separate Git branch. Because
the ports were fixed in advance, the branches merged without conflicts and the top
module could be wired directly. Table~\ref{tab:modules}
lists the modules.

\begin{figure}[!t]
\centering
\begin{tikzpicture}[font=\scriptsize, >={Stealth[length=1.5mm]},
  st/.style={draw, rounded corners=2pt, fill=white, minimum width=2.3cm, minimum height=0.62cm, align=center},
  node distance=0.32cm and 0.35cm]
\node[st, fill=stIF]  (a) {Shared header\\\texttt{defines.vh}};
\node[st, fill=stID, right=of a] (b) {Module RTL\\(per member branch)};
\node[st, fill=stEX, right=of b] (c) {Unit testbench\\PASS/FAIL};
\node[st, fill=stMEM, below=of c] (d) {Pull request\\and merge};
\node[st, fill=stWB, left=of d] (e) {Top-level\\integration};
\node[st, fill=hz, left=of e] (f) {System tests\\+ GTKWave};
\draw[->, thick] (a) -- (b);
\draw[->, thick] (b) -- (c);
\draw[->, thick] (c) -- (d);
\draw[->, thick] (d) -- (e);
\draw[->, thick] (e) -- (f);
\draw[->, dashed] (c.north) -- ++(0,0.25) -| node[pos=0.25, above, font=\tiny]{fix and re-run} (b.north);
\end{tikzpicture}
\caption{Interface-first design and verification flow used by the team.}
\label{fig:flow}
\end{figure}

\begin{table}[!t]
\centering
\caption{Verilog Modules of the Processor}
\label{tab:modules}
\setlength{\tabcolsep}{3pt}
\renewcommand{\arraystretch}{1.05}
\begin{tabular}{@{}lll@{}}
\toprule
\textbf{Module} & \textbf{Stage} & \textbf{Function} \\ \midrule
\texttt{pc\_reg}          & IF  & program counter with enable \\
\texttt{instr\_mem}       & IF  & 1024-word instruction ROM \\
\texttt{if\_id\_reg}      & IF/ID & latch, hold, flush to NOP \\
\texttt{reg\_file}        & ID  & 32$\times$32 registers, write-first \\
\texttt{imm\_gen}         & ID  & I/S/B/U/J immediates \\
\texttt{control\_unit}    & ID  & opcode decoder \\
\texttt{id\_ex\_reg}      & ID/EX & latch, flush to bubble \\
\texttt{alu\_control}     & EX  & selects 1 of 10 ALU ops \\
\texttt{alu}              & EX  & arithmetic, logic, shifts, compare \\
\texttt{branch\_unit}     & EX  & branch decision and target \\
\texttt{ex\_stage}        & EX  & forwarding muxes + above three \\
\texttt{ex\_mem\_reg}     & EX/MEM & pipeline register \\
\texttt{data\_mem}        & MEM & 1024-word data RAM \\
\texttt{mem\_wb\_reg}     & MEM/WB & pipeline register \\
\texttt{forwarding\_unit} & EX  & selects operand sources \\
\texttt{hazard\_unit}     & ID  & stall and flush control \\
\texttt{riscv\_pipeline\_top} & -- & integration, PC and WB muxes \\ \bottomrule
\end{tabular}
\end{table}

\subsection{Instruction Fetch (IF)}
The \texttt{pc\_reg} module holds the program counter. It loads
\texttt{pc\_next} on each clock edge when its enable \texttt{pc\_write} is high and holds
its value during a stall. The next-PC multiplexer selects
\begin{equation}
pc_{next} = \begin{cases} branch\_target, & pc\_src = 1\\ pc + 4, & \text{otherwise.}\end{cases}
\end{equation}
The \texttt{instr\_mem} module loads a program from a hex file with
\texttt{\$readmemh} and returns the word at \texttt{addr[11:2]}; unused words are
initialised to the canonical NOP (\texttt{addi x0,x0,0}). The IF/ID register latches the
PC and instruction; on reset or flush it outputs a NOP, and when \texttt{if\_id\_write}
is low it holds its contents.

\subsection{Instruction Decode (ID)}
The \texttt{control\_unit} decodes the 7-bit opcode into the control signals of
Table~\ref{tab:ctrl}. The \texttt{imm\_gen} module sign-extends the I, S, B, U and J
immediate formats according to the opcode. The \texttt{reg\_file} has 32 registers with
two combinational read ports and one write port; \texttt{x0} always reads as zero.
It implements a \emph{write-first bypass}: when WB writes register \texttt{rd} in the same
cycle that ID reads it, the new value is returned directly, which removes the WB-to-ID
hazard without any extra forwarding path. The ID/EX register passes all control and data
signals forward; on reset or flush it outputs all zeros, which is a bubble (no register
or memory write).

\begin{table}[!t]
\centering
\caption{Main Control Unit Truth Table}
\label{tab:ctrl}
\setlength{\tabcolsep}{2.6pt}
\renewcommand{\arraystretch}{1.1}
\begin{tabular}{@{}lccccccc@{}}
\toprule
\textbf{Opcode} & \textbf{RegW} & \textbf{MemR} & \textbf{MemW} & \textbf{M2R} & \textbf{ALUSrc} & \textbf{ALUOp} & \textbf{Other} \\ \midrule
R-type & 1 & 0 & 0 & 0 & 0 & 10 & -- \\
I-type & 1 & 0 & 0 & 0 & 1 & 11 & -- \\
LW     & 1 & 1 & 0 & 1 & 1 & 00 & -- \\
SW     & 0 & 0 & 1 & 0 & 1 & 00 & -- \\
Branch & 0 & 0 & 0 & 0 & 0 & 01 & branch \\
JAL    & 1 & 0 & 0 & 0 & 0 & 00 & jump \\
JALR   & 1 & 0 & 0 & 0 & 1 & 00 & jalr \\
LUI    & 1 & 0 & 0 & 0 & 1 & 00 & lui \\ \bottomrule
\end{tabular}
\end{table}

\subsection{Execute (EX)}
The \texttt{ex\_stage} wrapper contains the operand forwarding multiplexers, the
\texttt{alu\_control}, the \texttt{alu} and the \texttt{branch\_unit}. The ALU control maps
\texttt{alu\_op}, \texttt{funct3} and bit~5 of \texttt{funct7} to one of ten ALU
operations (ADD, SUB, AND, OR, XOR, SLL, SRL, SRA, SLT, SLTU). Operand~A is the forwarded
\texttt{rs1} value, or zero for LUI so that the ALU computes $0 + imm$. Operand~B is the
forwarded \texttt{rs2} value or the immediate, selected by \texttt{alu\_src}. The value
stored by SW is the forwarded \texttt{rs2}, so a store directly after the producing
instruction still writes the correct data.

The branch unit compares the forwarded operands for BEQ, BNE, BLT and BGE (signed) and
computes the target: $pc + imm$ for branches and JAL, and $(rs1 + imm)\,\&\,\sim\!1$ for
JALR. It asserts \texttt{pc\_src} when a branch is taken or for any jump. For JAL and
JALR the stage result is the link address:
\begin{equation}
ex\_result = (jump \lor jalr)\ ?\ pc + 4 : alu\_out .
\end{equation}

\subsection{Memory (MEM) and Write-Back (WB)}
The \texttt{data\_mem} module writes \texttt{wdata} on the rising edge when
\texttt{mem\_write} is set and returns \texttt{mem[addr[11:2]]} combinationally when
\texttt{mem\_read} is set (zero otherwise). The MEM/WB register carries the ALU result,
the loaded data, \texttt{rd} and the two remaining control bits. In WB the result is
selected as
\begin{equation}
wb\_data = mem\_to\_reg\ ?\ wb\_mem\_rdata : wb\_result
\end{equation}
and written to the register file when \texttt{wb\_reg\_write} is set.

\subsection{Hazard Handling}

\subsubsection{Structural hazards}
Separate instruction and data memories let IF and MEM access memory in the same cycle,
and the register file is read in ID and written in WB through separate ports. No
structural hazard remains.

\subsubsection{Data hazards and forwarding}
The \texttt{forwarding\_unit} compares the source registers of the instruction in EX with
the destination registers of the two older instructions:
\begin{align}
F_A &= \begin{cases}
10, & RegW_{M} \land rd_{M}\neq 0 \land rd_{M} = rs1_{E}\\
01, & RegW_{W} \land rd_{W}\neq 0 \land rd_{W} = rs1_{E}\\
00, & \text{otherwise}
\end{cases}
\end{align}
and $F_B$ is defined in the same way with $rs2_E$. Subscripts $M$ and $W$ denote the
EX/MEM and MEM/WB registers. The EX/MEM case is checked first because it holds the most
recent value of the register. Register \texttt{x0} is never forwarded.

\subsubsection{Load-use stall}
A loaded value exists only at the end of MEM, so it cannot be forwarded to an instruction
that enters EX in the very next cycle. The \texttt{hazard\_unit} detects this case in ID:
\begin{equation}
load\_use = MemR_E \land rd_E \neq 0 \land (rd_E = rs1_D \lor rd_E = rs2_D).
\end{equation}
When it is true, \texttt{pc\_write} and \texttt{if\_id\_write} are cleared, so the PC and
IF/ID hold, and \texttt{id\_ex\_flush} inserts a bubble. In the next cycle the load is in
WB and its value is forwarded from MEM/WB. Fig.~\ref{fig:timing} and
Table~\ref{tab:loaduse} show this for the program
\texttt{lw x2,0(x0); add x3,x2,x2}. The table was taken from the simulation trace of
system test~4.

% ---------------- Pipeline timing figure ----------------
\begin{figure}[!t]
\centering
\begin{tikzpicture}[x=0.78cm, y=0.62cm, font=\scriptsize]
\tikzset{
  c/.style={draw, minimum width=0.74cm, minimum height=0.5cm, inner sep=0pt, font=\scriptsize\bfseries},
}
\foreach \i in {1,...,8} \node at (\i,0.75) {C\i};
\node[anchor=east] at (0.5,0) {\texttt{lw x2,0(x0)}};
\node[anchor=east] at (0.5,-1) {\texttt{add x3,x2,x2}};
\node[anchor=east] at (0.5,-2) {\texttt{jal x0,0}};
% lw
\node[c, fill=stIF]  at (1,0) {IF};
\node[c, fill=stID]  at (2,0) {ID};
\node[c, fill=stEX]  at (3,0) {EX};
\node[c, fill=stMEM] at (4,0) {MEM};
\node[c, fill=stWB]  at (5,0) {WB};
% add
\node[c, fill=stIF]  at (2,-1) {IF};
\node[c, fill=stID]  at (3,-1) {ID};
\node[c, fill=hz]    at (4,-1) {ID};
\node[c, fill=stEX]  at (5,-1) {EX};
\node[c, fill=stMEM] at (6,-1) {MEM};
\node[c, fill=stWB]  at (7,-1) {WB};
% jal
\node[c, fill=stIF]  at (3,-2) {IF};
\node[c, fill=hz]    at (4,-2) {IF};
\node[c, fill=stID]  at (5,-2) {ID};
\node[c, fill=stEX]  at (6,-2) {EX};
\node[c, fill=stMEM] at (7,-2) {MEM};
\node[c, fill=stWB]  at (8,-2) {WB};
% bubble
\node[c, dashed, fill=gray!15, font=\scriptsize] at (4,-3) {bub.};
\node[anchor=east] at (0.5,-3) {\textit{bubble in EX}};
% forwarding arrow
\draw[-{Stealth[length=1.6mm]}, very thick, blue!70!black] (5,-0.25) .. controls (5.3,-0.5) .. (5,-0.75);
\node[blue!70!black, anchor=west, font=\tiny] at (5.55,0) {forward MEM/WB $\rightarrow$ EX};
\end{tikzpicture}
\caption{Pipeline timing for a load-use dependency. The dependent \texttt{add} is held in
ID for one cycle (yellow) while a bubble enters EX; the loaded value is then forwarded from
MEM/WB.}
\label{fig:timing}
\end{figure}

\begin{table}[!t]
\centering
\caption{Simulated Hazard Signals for the Load-Use Example (Test 4)}
\label{tab:loaduse}
\setlength{\tabcolsep}{3pt}
\renewcommand{\arraystretch}{1.1}
\begin{tabular}{@{}cccccc@{}}
\toprule
\textbf{Cycle} & \textbf{PC} & \textbf{Instr. in ID} & \textbf{pc\_write} & \textbf{id\_ex\_flush} & $F_A/F_B$ \\ \midrule
4 & 0x0C & \texttt{lw x2}          & 1 & 0 & 00/10 \\
5 & 0x10 & \texttt{add x3,x2,x2}   & \textbf{0} & \textbf{1} & 00/00 \\
6 & 0x10 & \texttt{add x3,x2,x2}   & 1 & 0 & 00/00 \\
7 & 0x14 & \texttt{jal x0,0}       & 1 & 0 & \textbf{01/01} \\ \bottomrule
\end{tabular}
\end{table}

\subsubsection{Control hazards}
Branches and jumps are resolved in EX with a predict-not-taken policy. When
\texttt{pc\_src} is asserted, the two younger instructions in IF/ID and ID/EX are on the
wrong path, so the hazard unit sets \texttt{if\_id\_flush} and \texttt{id\_ex\_flush}
and loads the target into the PC. A taken branch therefore costs two cycles; a not-taken
branch costs nothing. Flush has priority over a simultaneous load-use stall, because the
stalled instruction is itself on the wrong path. The complete hazard unit is:
\begin{center}
\footnotesize\setlength{\tabcolsep}{3pt}
\begin{tabular}{@{}lcccc@{}}
\toprule
\textbf{Condition} & \textbf{pc\_write} & \textbf{if\_id\_write} & \textbf{if\_id\_flush} & \textbf{id\_ex\_flush} \\ \midrule
\texttt{pc\_src}          & 1 & 1 & 1 & 1 \\
\texttt{load\_use} only   & 0 & 0 & 0 & 1 \\
neither                   & 1 & 1 & 0 & 0 \\ \bottomrule
\end{tabular}
\end{center}
Fig.~\ref{fig:hazflow} shows the same decision order as a flowchart.
\begin{figure}[!t]
\centering
\begin{tikzpicture}[font=\scriptsize, >={Stealth[length=1.5mm]},
  dec/.style={draw, diamond, aspect=2.4, fill=hz, inner sep=1pt, align=center},
  act/.style={draw, rounded corners=2pt, fill=white, align=center, minimum width=2.6cm, minimum height=0.55cm}]
\node[draw, rounded corners=6pt, fill=gray!15] (s) at (0,0) {each cycle};
\node[dec] (d1) at (0,-1.05) {\texttt{pc\_src}?};
\node[act, fill=stWB] (a1) at (3.0,-1.05) {flush IF/ID and ID/EX\\PC $\leftarrow$ target};
\node[dec] (d2) at (0,-2.5) {\texttt{load\_use}?};
\node[act, fill=stEX] (a2) at (3.0,-2.5) {hold PC and IF/ID\\bubble into ID/EX};
\node[act, fill=stID] (a3) at (0,-3.75) {normal advance\\forwarding unit selects operands};
\draw[->] (s) -- (d1);
\draw[->] (d1) -- node[above]{yes} (a1);
\draw[->] (d1) -- node[right]{no} (d2);
\draw[->] (d2) -- node[above]{yes} (a2);
\draw[->] (d2) -- node[right]{no} (a3);
\end{tikzpicture}
\caption{Decision order of the hazard detection unit. A taken branch or jump is checked
first, so flushing always overrides a load-use stall.}
\label{fig:hazflow}
\end{figure}

Fig.~\ref{fig:flush} shows the flush observed for \texttt{beq x1,x2,8} in system test~5:
the branch is taken in cycle~5 and the instructions at 0x0C and 0x10 are discarded.

\begin{figure}[!t]
\centering
\begin{tikzpicture}[x=0.78cm, y=0.62cm, font=\scriptsize]
\tikzset{c/.style={draw, minimum width=0.74cm, minimum height=0.5cm, inner sep=0pt, font=\scriptsize\bfseries}}
\foreach \i in {3,...,9} \node at (\i-2,0.75) {C\i};
\node[anchor=east] at (0.5,0)  {\texttt{0x08 beq x1,x2,8}};
\node[anchor=east] at (0.5,-1) {\texttt{0x0C addi x3,..}};
\node[anchor=east] at (0.5,-2) {\texttt{0x10} (wrong path)};
\node[anchor=east] at (0.5,-3) {\texttt{0x10} (target)};
\node[c, fill=stIF] at (1,0) {IF};
\node[c, fill=stID] at (2,0) {ID};
\node[c, fill=stEX, line width=1.2pt, draw=red!70!black] at (3,0) {EX};
\node[c, fill=stMEM] at (4,0) {MEM};
\node[c, fill=stWB] at (5,0) {WB};
\node[c, fill=stIF] at (2,-1) {IF};
\node[c, fill=stID] at (3,-1) {ID};
\node[c, dashed, fill=gray!15] at (4,-1) {--};
\node[c, fill=stIF] at (3,-2) {IF};
\node[c, dashed, fill=gray!15] at (4,-2) {--};
\node[c, fill=stIF] at (4,-3) {IF};
\node[c, fill=stID] at (5,-3) {ID};
\node[c, fill=stEX] at (6,-3) {EX};
\node[c, fill=stMEM] at (7,-3) {MEM};
\draw[red!70!black, thick] (3.45,-0.85) -- (3.45,-2.15);
\node[red!70!black, anchor=west, font=\tiny, align=left] at (4.45,-1.5) {flush IF/ID,\\ ID/EX};
\end{tikzpicture}
\caption{Branch resolution in EX (cycle 5) and flushing of the two wrongly fetched
instructions, as observed in system test 5. The taken branch costs two cycles.}
\label{fig:flush}
\end{figure}

\subsection{Implementation and Tools}
The design consists of 17 Verilog modules sharing one header (\texttt{defines.vh}) for
opcodes, ALU control codes and the NOP encoding~\cite{ieee1364}. Port names follow a
fixed rule: a pipeline-register input carries its source-stage prefix and the matching
output carries the destination-stage prefix (for example \texttt{ex\_result} to
\texttt{mem\_result}). The top module \texttt{riscv\_pipeline\_top} instantiates and wires
all modules and contains the PC and write-back multiplexers. Simulation uses Icarus
Verilog~\cite{iverilog} and waveforms are viewed in GTKWave~\cite{gtkwave}. Source code is
version-controlled on GitHub, with each member working on a separate branch merged through
pull requests.

% ---------------------------------------------------------------------------
